\section{Matching the edge potentials at the vertices}
\label{sec:vertices}

We now pass from the cut manifolds to the exact algebraic data of
Theorem~\ref{thm:tensor-obstruction}.  All residual ideals in this section are
ideals of formal rings, without passage to their radicals.

\begin{proposition}[Vertex data]
\label{prop:vertex-data}
Suppose that the data of Proposition~\ref{prop:geometric-reduction} exist,
with $m\geq 5$, $p\geq9$ extra letters, and the finite family of
independently conjugated boundary laws described there.  Their Brunnian
property means that omitting any slot makes the law the identity, as in
Equation~\eqref{eq:grope-word-properties}.
There are a characteristic-zero field $K$, internal coordinate blocks $h_i$
for the $2m$ edges, and algebraic power-series germs $F_i(h_i)$ and $b_i(h_i)$,
the latter for the $m$ active edges, with the following properties.
Put
\[
 T=K[t_i:i\text{ active}]/(t_i^2:i\text{ active}),\qquad
 S_0=\sum_i F_i,
 \qquad S_t=S_0+\sum_{i\text{ active}}t_i b_i.
\]
There are formal parametrizations $H_L(x,t)$ over $T$ and $H_R(y)$ over $K$
whose central tangent maps are injective and have complementary images in
the full internal coordinate space.  They are smooth graph germs over $T$
and $K$, respectively, and
\begin{equation}
\label{eq:vertex-final-vanishing}
 S_t(H_L)=0,\qquad S_0(H_R)=0.
\end{equation}
For every three distinct active indices $i,j,k$,
\begin{align}
 t_i t_j t_k
 &\in\left((\partial_{h_a}S_t)(H_L):a=1,\ldots,N\right)
       \subset T[[x]],                                      \label{eq:vertex-final-source}\\
 (b_i b_j b_k)(H_R)
 &\in\left((\partial_{h_a}S_0)(H_R):a=1,\ldots,N\right)
       \subset K[[y]],                                      \label{eq:vertex-final-target}
\end{align}
where $N$ is the total number of internal coordinates.  The field can be
taken to be $K=\mathbb C((\sigma))$.  The individual functions can be
normalized so that $F_i(0)=b_i(0)=0$.
\end{proposition}

We first construct formal sections from the two vertex connection spaces.
Their curvature ideals will remain fixed as we normalize the boundary
restrictions and correct the sections to make the sums of potentials
vanish.  We then pass to algebraic edge coordinates and use the boundary
laws to obtain the triple memberships.  The final step replaces the
logarithmic perturbations by algebraic functions and corrects the source
section once more.

Orient each $Y_i$ as a boundary face of $A$.  Its orientation as a face
of $B$ is the opposite one.  Write
$\alpha_i$ for its surviving peripheral direction and $\beta_i$ for the
other direction.  The plus endpoint is the one at which $\beta_i$ bounds
in the free boundary piece; the minus endpoint is the one at which
$\alpha_i$ bounds.  An edge is \emph{active} when its plus endpoint is at
$A$.  There are $m$ active edges.

Lemma~\ref{lem:connected-tree-cuts} gives smooth cuts that miss the locally
flat disks.  Their complementary pieces are smooth manifolds with corners,
so ordinary smooth forms and Stokes' formula apply, with matching
pullbacks on their smooth faces at corners.

\subsection{The linear restriction maps}

Put $\g=\mathfrak{sl}_2(\mathbb C)$ and use the invariant pairing
$\kappa(U,V)=\tr(UV)$.  The internal edge space is
\[
 \mathcal H=\bigoplus_i
 H^1(Y_i,\partial Y_i;\mathbb C)\otimes\g.
\]
For a connected oriented three-manifold with connected torus boundary,
the natural map from relative degree-one cohomology identifies it with
the kernel of restriction to the boundary.  Thus these are precisely the
internal spaces in Proposition~\ref{prop:edge-potential}.

At this linear stage, fixing an evaluation means requiring its tangent
variation to vanish.  For \(P=A,B\), let \(\mathcal V_P\) be the common
kernel in \(H^1(P;\mathbb C)\otimes\g\) of evaluation on the extra
letters at \(P\) and on the \(\alpha_i\) whose plus endpoint is \(P\).
The degree-one assertion below is the isomorphism
\(\mathcal V_A\oplus\mathcal V_B\to\mathcal H\).  Later prescribed
holonomies need not be the identity; it is their variations that are
fixed in this tangent-space statement.

\begin{lemma}
\label{lem:vertex-linear}
After fixing the evaluations on the extra letters and on all plus
meridians, the direct sum of the two vertex degree-one spaces maps
isomorphically, by difference of restrictions, to $\mathcal H$.
Moreover,
\begin{equation}
\label{eq:vertex-h2-isomorphism}
 H^2(A;\mathbb C)\oplus H^2(B;\mathbb C)
       \longrightarrow\bigoplus_i H^2(Y_i;\mathbb C)
\end{equation}
is an isomorphism.  In particular each vertex summand maps injectively.
\end{lemma}

\begin{proof}
Lemma~\ref{lem:geometric-edge-cohomology} gives
Equation~\eqref{eq:vertex-h2-isomorphism} and the degree-one restriction
isomorphism after the extra evaluations are fixed.

For each edge the image of $H^1(Y_i)$ in the cohomology of its boundary
torus is a line.  Restrictions from the two vertex sides span $H^1(Y_i)$,
and their torus restrictions lie on the two coordinate axes, because the
two respective longitudes bound in the boundary pieces.  This line must
be one coordinate axis.  Evaluation on its surviving loop $\alpha_i$
comes only from the plus vertex.  Consequently, after the extra
evaluations have been fixed, the inverse image of
$\bigoplus_i H^1(Y_i,\partial Y_i)$ is exactly the subspace on which all
plus evaluations vanish.  The asserted restricted map is therefore an
isomorphism.  Tensoring these statements with $\g$ proves the result.
\end{proof}

The pairing
\begin{equation}
\label{eq:vertex-perfect-edge-pairing}
 \bigoplus_i H^1(Y_i,\partial Y_i)\otimes\g
 \ \times\ \bigoplus_i H^2(Y_i)\otimes\g
 \longrightarrow\mathbb C,
 \qquad (u,v)\longmapsto\sum_i\int_{Y_i}\kappa(u\wedge v)
\end{equation}
is perfect by Poincar\'e--Lefschetz duality.  We will use both the
degree-one isomorphism and this degree-two pairing.

\subsection{Marked formal backgrounds and vertex slices}

The boundary laws will be tested at a finite family of conjugating
matrices whose adjoint operators span every endomorphism of $\g$.
Introduce a formal parameter $\sigma$ and set $K=\mathbb C((\sigma))$.
Choose nine matrices in $SL_2(\mathbb C[[\sigma]])$, one the identity,
all equal to the identity at $\sigma=0$, whose adjoint operators span
$\End_K(\g\otimes K)$. Equation~\eqref{eq:vertex-explicit-background}
gives an explicit family, and Lemma~\ref{lem:vertex-adjoint-span} verifies
its spanning property. The family is fixed independently of the cuts.
Its independent actions on the three slots will detect every triple
coefficient in the boundary laws. Assign the
eight nonidentity values to eight extra letters at each vertex and set
the remaining extra holonomies to the identity.  Write
$g_s(\sigma)\in SL_2(\mathbb C[[\sigma]])$, with $g_s(0)=1$, for these
assigned holonomies.  These choices extend to a flat background
on the whole ambient bipartite model with trivial meridian holonomies.
To see the marking explicitly, choose the paths through the joins as
edge paths of the free graph groupoid.  The extra loops are its vertex
generators, and edge transports can be chosen freely.  Choose those
transports so that the two torus frames at every join agree with the
frames determined by the boundary basing paths.  This simultaneously
fixes the markings used for holonomy and for the boundary laws.

Initially all constructions take place jointly formally at
$\sigma=0$.  The restricted flat backgrounds have coordinates of positive
$\sigma$ order in the central vertex and edge charts.  We do not yet
invert $\sigma$.

Fix $P=A$ or $P=B$, with its basepoint in $K_P$.  In degree zero use smooth
$\g$-valued functions vanishing at that point, and in positive degrees
use unrestricted forms on $P$.  Choose a contraction onto cohomology,
denoted by inclusion $\iota$, projection $\pi$ and homotopy $\mathsf K_P$, with
\[
 \dd\mathsf K_P+\mathsf K_P\dd=1-\iota\pi,
 \qquad \mathsf K_P^2=\mathsf K_P\iota=\pi\mathsf K_P=0.
\]
Degree-zero cohomology is zero.  For
$y_{\rm full}\in H^1(P;\mathbb C)\otimes\g$, the recursion
\[
 a_P=\iota y_{\rm full}-\tfrac12\mathsf K_P[a_P,a_P]
\]
defines a unique formal one-form.  Write
\[
 \mathcal F_P=\dd a_P+\tfrac12[a_P,a_P],\qquad r_P=\pi\mathcal F_P.
\]
The list $r_P$ has $q_P=\dim(H^2(P;\mathbb C)\otimes\g)$ entries, of
order at least two.  The slice equation and Bianchi identity give
\begin{equation}
\label{eq:vertex-residual-curvature}
 \mathcal F_P=\iota r_P-\mathsf K_P[a_P,\mathcal F_P]
       =(1+\mathsf K_P\ad_{a_P})^{-1}\iota r_P.
\end{equation}
In particular the curvature ideal is exactly the ideal generated by
$r_P$, and its leading coefficient forms are the selected representatives
of $H^2(P)\otimes\g$.  Successively removing exact one-form coordinates
by based gauges shows that every based flat formal connection is
represented by this slice modulo $r_P$.  These constructions are
functorial over nilpotent coefficient rings and pass to their inverse
limits.  No regular-sequence assertion about $r_P$ is made.

Fix nonzero $T_i\in\g$ for the active edges.  We use the connection
convention of Proposition~\ref{prop:edge-potential}, in which the strip
connection $u\zeta$ has holonomy $\exp(-u)$.  Impose on the vertex slice
the exact logarithmic holonomies
\begin{itemize}
\item $\log g_s(\sigma)$ on the extra letters;
\item $-t_iT_i$ on an active plus meridian at $A$;
\item zero on every plus meridian at $B$.
\end{itemize}
There are no parameters on the other plus ends.  Logarithmic holonomy
along a fixed based loop is a formal function of the off-shell
connection, whose linear term is minus integration of the connection along
that loop.  Lemma~\ref{lem:vertex-linear} therefore makes these conditions
a smooth formal parameter slice.  Keep its other coordinates $y_P$
free.  The resulting coefficient rings are
\[
 R_A=\mathbb C[[\sigma,y_A]][t_i:i\text{ active}]/(t_i^2),
 \qquad R_B=\mathbb C[[\sigma,y_B]],
\]
completed at their displayed joint origins.  Write $I_P=(r_P)$ for the
substituted residual ideal.  It is contained in the square of the joint
maximal ideal.  Modulo $I_P$ there is an actual based flat system with
exactly the prescribed extra and plus holonomies.  We keep this list
$r_P$ as functions on the parameter domain throughout the subsequent
normalizations and section corrections.  Thus the flat quotient
$R_P/I_P$ stays fixed even when its chosen lifts to forms and edge
coordinates change.

For an active edge, the induced holonomies in the flat quotient
$R_A/I_A$ satisfy
\[
 G_{\alpha_i}=\exp(-t_iT_i),\qquad G_{\beta_i}=1.
\]
In the flat quotient $R_B/I_B$ they satisfy $G_{\alpha_i}=1$, while the
marked $G_{\beta_i}$ varies.  The central edge potential there contributes
$-f_i$ because the boundary orientation is reversed, where $f_i$ is the
central potential of Proposition~\ref{prop:edge-potential}.  These
holonomy identities concern the respective flat quotients, not the
off-shell lifts.

\subsection{Boundary normalization and lifts of the flat quotient}

The quotient $R_P/I_P$ carries the flat systems and their group laws.
The tensor theorem, however, requires sections in the ambient edge
coordinates, where the curvature need not vanish.  We first normalize
the restrictions on the flat quotient and then lift them to these
ambient, or off-shell, coordinates.  Changes in the lifts will be
measured against the fixed list $r_P$.

Recall that $K_P$ is a connected sum of $p$ copies of $S^1\times S^2$
with tubes about an unlink removed.  In the unlink exterior choose
disjoint proper disks dual to the port meridians, with boundaries the
unlink longitudes.  In the connected-sum pieces choose disjoint
nonseparating spheres dual to the extra $S^1$ generators.  Thin mutually
disjoint collars of these disks and spheres carry closed Thom
one-forms $\rho_\nu$, normalized to evaluate to one on the corresponding
generator.  The disk strips meet the port tori in the coordinate
densities dual to the port meridians; the sphere strips miss those tori.
Choose the basing arcs off the strips.  Since the port meridians and
extra letters freely generate $\pi_1K_P$, a flat system with their
specified holonomies has the representative
\[
 a_{K_P}=\sum_\nu M_\nu\rho_\nu,
\]
where the $M_\nu$ are the corresponding signed logarithms.  Each strip
form is closed and their cross wedges vanish, so this connection is
flat for arbitrary matrix coefficients.
At a plus torus its trace is $u_i\zeta_i$, and at a minus torus its trace
is a multiple of $\xi_i$.  The two strip choices are the compatible
torus choices of Proposition~\ref{prop:edge-potential}.  The area
representatives can be supported away from both strips.  In particular,
\begin{equation}
\label{eq:vertex-k-cs-zero}
 \mathcal C_{K_P}(a_{K_P})=0.
\end{equation}

Normalize first over $R_P/I_P$.  Based parallel transport gives a gauge
from the restriction of $a_P$ to this strip representative.  After this gauge,
the flat restriction on every $Y_i$ can be put into the edge slice
$s_i(h_i,t_i)$, or $s_i(h_i,0)$ as appropriate, by a gauge equal to the
identity on its torus.  This gives coordinates $h_i^P(y_P)$ over
$R_P/I_P$.  The normalizations are compatible on intersections of
boundary faces: the second gauges are the identity on the tori.

To lift these quotient-ring data to actual formal forms, choose a
power-series presentation of the Noetherian complete ring $R_P$ and
fix continuous linear division operators
for the finite list $r_P$ and for passage to $R_P/I_P$.  Such operators
are obtained by formal division, or by choosing compatible filtered
linear splittings.  Applied coefficient by coefficient, they lift
scalar series and write a form-valued discrepancy in $I_P$ as
$\sum_a r_{P,a}\gamma_a$ with form-valued series $\gamma_a$.
Each output coefficient involves only finitely many input coefficients.
Thus the same operators apply to the vector spaces of smooth functions
and differential forms.  This is the actual finite-list ideal
membership used below.

Lift the near-identity gauges by first lifting their Lie-algebra-valued
logarithms.  Smooth collar extension extends each logarithm off the
relevant boundary face; the compatibility at the tori allows the
extensions to be combined.  Lift the coordinates $h_i^P$ as well.
Finally change the connection by $I_P$-multiples of one-forms, extended
over a collar of $K_P$, so that its restriction to $K_P$ is
\emph{exactly} the strip representative with the constrained plus
coefficients.  Denote this trial connection by $\widehat a_P$.  It has
the following properties:
\begin{enumerate}
\item its curvature belongs to $I_P$;
\item its restriction to $Y_i$ equals the specified edge slice modulo
      $I_P$;
\item to first order in the actual residual list, modulo the joint
      maximal ideal times that list, its curvature representatives are
      those in Equation~\eqref{eq:vertex-residual-curvature}, changed only by
      exact two-forms.
\end{enumerate}
For the last statement, write an added one-form as
$\eta=\sum_a r_{P,a}\gamma_a$ and put
$\mathcal F(a)=\dd a+\tfrac12[a,a]$ for an intermediate connection $a$.
The source differential treats the
functions $r_{P,a}$ of the formal parameters as constants.  Hence
\[
\begin{aligned}
 \mathcal F(a+\eta)
 &=\mathcal F(a)+\sum_a r_{P,a}\,\dd\gamma_a
                  +[a,\eta]+\tfrac12[\eta,\eta]\\
 &\equiv\mathcal F(a)+\sum_a r_{P,a}\,\dd\gamma_a
                  \pmod{\mathfrak m_P I_P}.
\end{aligned}
\]
The congruence is coefficientwise, and $\mathfrak m_P$ is the joint
maximal ideal: $a$ has positive
order and $I_P\subset\mathfrak m_P^2$.  A gauge with identity value
at the joint origin also changes curvature only modulo
$\mathfrak m_PI_P$.  Thus the leading coefficient forms of the fixed
residual list change by exact two-forms.  This calculation uses a
chosen expression for $\eta$ and requires no independence of the
$r_{P,a}$.  The lifts have the same first derivatives as the original
quotient coordinates because $I_P$ has order at least two.

\subsection{The square-ideal and gradient identities}

We have lifted the flat boundary restrictions to ambient edge
coordinates without changing them modulo the fixed residual ideal.
The next calculation measures the error in the sum of their
potentials and identifies the restricted gradient ideal. These two
statements will permit correction of the sections while preserving
their flat-system markings.

Let $f_i$ and $b_i^{\rm form}$ be the edge functions before algebraizing
the perturbations.  With the orientations already chosen, put
\[
 S_A=\sum_i f_i+\sum_{i\text{ active}}t_i b_i^{\rm form},
 \qquad S_B=-\sum_i f_i.
\]
Write $H_P$ for the lifted parametrization consisting of all $h_i^P$.

\begin{lemma}
\label{lem:vertex-offshell-identities}
At the joint origin the tangent images of $H_A$ and $H_B$, with the
external parameters fixed, are complementary.  For each vertex,
\begin{equation}
\label{eq:vertex-square-and-gradient}
 S_P(H_P)\in I_P^2,
 \qquad (\partial_h S_P)(H_P)=M_P r_P,
\end{equation}
where $M_P$ is an actual coefficient matrix whose reduction at the
joint origin has full column rank $q_P$.  Consequently the restricted
gradient ideal equals $I_P$.
\end{lemma}

\begin{proof}
The assertion about tangent images is Lemma~\ref{lem:vertex-linear}:
the linearization of restriction to the edge slices is the restriction
map into the relative degree-one edge spaces.  Thus each parametrization
is an immersion and can be written as a graph after a linear projection.

For the first identity, Stokes' formula gives
\[
 \mathcal C_{\partial P}(\widehat a_P|_{\partial P})
       =\tfrac12\int_P\tr(\mathcal F_{\widehat a_P}
                              \wedge\mathcal F_{\widehat a_P})
       \in I_P^2.
\]
The $K_P$ contribution is zero by Equation~\eqref{eq:vertex-k-cs-zero}.  On an
edge interpolate between the trial restriction and its chosen slice.
Their difference is in $I_P$, and every curvature along the interpolation
is in $I_P$.  The interior variation of the action is therefore in
$I_P^2$.  On a minus end both boundary traces lie in the $\xi_i$
polarization, so the boundary variation vanishes.  The same holds at a
frozen plus end, where $u_i=0$.  At an active plus end the trial trace
is $u_i\zeta_i$, the slice trace is $u_i\zeta_i+v_i\xi_i$, and the
boundary variation is
\[
 -\tfrac{\epsilon_i}{2}\tr(u_i v_i).
\]
It is canceled exactly by the adjustment in
Proposition~\ref{prop:edge-potential}.  This proves $S_P(H_P)\in I_P^2$
with an actual quadratic expression in the residual list.

The fixed-$u$ variation formula in Equation~\eqref{eq:edge-fixed-u-gradient} expresses the
internal gradient entries as integrations of the slice curvatures
against derivatives of the slices.  Since those curvatures vanish
modulo $I_P$, all entries belong to $I_P$.  To determine a coefficient
matrix, start with Equation~\eqref{eq:vertex-residual-curvature} and the explicit
$I_P$-multiple differences just constructed.  At the joint origin, the
coefficient of $r_{P,a}$ in the $j$th gradient entry is, with its
boundary-orientation sign,
\begin{equation}
\label{eq:vertex-leading-matrix}
 \int_{Y_i}\kappa\bigl(\theta_{i,j}\wedge
                     \omega_{P,a}|_{Y_i}\bigr).
\end{equation}
Here $\theta_{i,j}$ is the closed relative one-form representing the
$j$th edge tangent and $\omega_{P,a}$ is the chosen closed two-form
representative of the residual coordinate.  Added leading exact
two-forms do not change this number: integrate by parts and use the
zero boundary pullback of $\theta_{i,j}$.  By
Equation~\eqref{eq:vertex-h2-isomorphism} and the perfect pairing in
Equation~\eqref{eq:vertex-perfect-edge-pairing}, the matrix in
Equation~\eqref{eq:vertex-leading-matrix} has full column rank.

Choose such a maximal minor.  It is a unit in $R_P$, so the corresponding
gradient entries recover the list $r_P$ by its inverse.  This proves
equality of ideals.  We used the matrix obtained from the curvature
expansion, not an inference that the residuals are independent modulo
some ideal.  Possible syzygies among them have no effect on the argument.
\end{proof}

We record the formal correction argument used at the joint origin.

\begin{lemma}[Corrections with a fixed residual list]
\label{lem:vertex-correction}
Let $R$ be a complete Noetherian local $\mathbb Q$-algebra with maximal
ideal $\mathfrak m_R$, let $I=(r_1,\ldots,r_q)$, let
$S\in R[[h_1,\ldots,h_N]]$, and let $H\in\mathfrak m_R^N$.
Internal derivatives are taken over $R$, and substitution at $H$ is
$\mathfrak m_R$-adically continuous.  Suppose
$S(H)\in I^2$ and $(\partial_h S)(H)=Mr$, with a unit full-column minor
in $M$.  Assume also that $I\subset\mathfrak m_R^2$ and the internal
Hessian of $S$ evaluated at $H$ has entries in $\mathfrak m_R$.
Then there is a convergent correction of $H$ by $I$-multiples making
$S(H)=0$.  It is unchanged modulo $I$ and to first order at the origin;
the gradient still generates $I$.  The functions $r_a$ on the parameter
domain are held fixed throughout the correction.
\end{lemma}

\begin{proof}
If the current error is $E=\sum_{a,b}c_{ab}r_ar_b$, use the inverse
minor of $M$ to choose a vector $\delta H$ of $I$-multiples with
$(\partial_h S)(H)\cdot\delta H=-E$.  More generally, if the coefficients
$c_{ab}$ lie in an ideal $Q$, choose $\delta H\in QI$.  This is a linear
calculation using the displayed list and its matrix; it does not require
unique expressions for elements of $I$ or $I^2$.

An error in $\mathfrak m_R^n I^2$ gives a
correction in $\mathfrak m_R^n I$.  The quadratic Taylor term is
$\frac12\delta H^{\mathsf T}(\operatorname{Hess}S)(H)\delta H$;
the factor $1/2$ and all higher Taylor coefficients are defined because
$R$ is a $\mathbb Q$-algebra.  This term belongs to
$\mathfrak m_R^{2n+1}I^2$ by the Hessian hypothesis, and higher terms
belong to the same ideal because $I\subset\mathfrak m_R^2$.
At the $k$th step write
$\delta H_k=D_kr$, with every entry of $D_k$ in
$\mathfrak m_R^{n_k}$, where $n_0=0$ and $n_{k+1}=2n_k+1$.
The Hessian remains in $\mathfrak m_R$, so the gradient coefficient
matrix changes by a matrix in $\mathfrak m_R^{n_k+1}$.
Thus both $\sum_kD_k$ and the sequence of gradient coefficient matrices
converge.  The total correction is $(\sum_kD_k)r\in I$, and the limiting
gradient coefficient matrix has the same reduction as the original
one, hence retains the chosen unit full-column minor.  Continuity now
gives $S(H_\infty)=0$ and $(\partial_h S)(H_\infty)=M_\infty r$.  The
assumption $I\subset\mathfrak m_R^2$ preserves the tangent map.

\end{proof}

Apply Lemma~\ref{lem:vertex-correction} to the identities of
Lemma~\ref{lem:vertex-offshell-identities}.
The central edge functions have order at least three; the other terms
carry a parameter $t_i$.  Hence the internal Hessian vanishes at the
joint origin.  We obtain corrected sections, still denoted by $H_A,H_B$,
with
\begin{equation}
\label{eq:vertex-exact-formal-sections}
 S_A(H_A)=S_B(H_B)=0,
 \qquad ((\partial_h S_P)(H_P))=I_P.
\end{equation}
They retain their complementary tangents and their values modulo $I_P$.
In particular they retain exactly the marked edge holonomies on the
flat residual schemes.

\subsection{Algebraic central coordinates and the order of completion}
\label{rem:tensor-recentering}

The corrected sections already satisfy exact vanishing and retain the
marked holonomies modulo their residual ideals. We now make the
central edge potentials algebraic and move the common background to
the origin. The order of this recentering and the subsequent passage
to Laurent coefficients is part of the construction.

Apply Theorem~\ref{thm:edge-algebraicity} separately on the edges,
initially at the trivial system.  Thus, after separate invertible formal
changes of the $h_i$, the central functions are algebraic germs $F_i$.
On their critical schemes the marked $\beta_i$ holonomies have algebraic
lifts $G_i$ in the smooth ambient germs.  Choose these lifts with
$G_i(0)=1$.  The functions $b_i^{\rm form}$ satisfy
\begin{equation}
\label{eq:vertex-log-mod-jacobian}
 b_i^{\rm form}\equiv -\epsilon_i\tr(T_i\log G_i)
                                      \pmod{\Jac(F_i)}.
\end{equation}
Here the sign is the one fixed by $\operatorname{Hol}_{\beta_i}=\exp(-v_i)$.
Write the Jacobian-ideal difference as
$\sum_a c_{i,a}\partial_{h_{i,a}}F_i$.  A separate coordinate change
linear in $t_i$ absorbs it: the identity
\[
 F_i(h_i+t_i c_i(h_i))
     =F_i(h_i)+t_i\sum_a c_{i,a}(h_i)\partial_{h_{i,a}}F_i(h_i)
\]
is exact because $t_i^2=0$.  Choosing the sign according to the direction
of the change replaces the coefficient by
\[
 \ell_i(G_i)=-\epsilon_i\tr(T_i\log G_i).
\]
The changes are separate in the edge blocks and preserve all equations
and ideal memberships on the transformed sections.  At parameter zero
the changes linear in $t_i$ are the identity.

Let $y_P^0(\sigma)$ and $h_i^0(\sigma)$ be the coordinates of the fixed
flat background in these charts.  They have positive $\sigma$ order.
Both vertex sections contain the same edge background, because all
earlier corrections vanish on their residual schemes.  Recenter by
$y_P=y_P^0(\sigma)+z_P$ and $h_i=h_i^0(\sigma)+z_i$ \emph{before}
passing to Laurent coefficients.  Substitution gives continuous maps
\[
 \mathbb C[[\sigma,y_P]]\longrightarrow\mathbb C[[\sigma,z_P]],
 \qquad
 \mathbb C[[\sigma,h_i]]\longrightarrow\mathbb C[[\sigma,z_i]].
\]
For each coefficient in the new internal variables, the defining sum
converges $\sigma$-adically.  We then include the coefficient ring in
$K=\mathbb C((\sigma))$ and complete in the new internal variables.
We rename them $y_P,h_i$.

This order also describes the changes already used: if $\phi$ is a
separate original formal coordinate change, its recentered expression is
\[
 z\longmapsto\phi(h_i^0(\sigma)+z)-\phi(h_i^0(\sigma)),
\]
formed over $\mathbb C[[\sigma,z]]$ first.  Later inverses over $K$ may
have coefficients with unbounded negative $\sigma$ orders.  They are
only composed with centered series, or with series whose constant terms
are nilpotent parameter multiples; these substitutions are well defined
in $K[[z]][t_i]/(t_i^2)$.  In particular, an absorbed Jacobian term gives
a separate square-zero change
\[
 z_i\longmapsto z_i+t_i c_i(z_i),\qquad
 z_i\longmapsto z_i-t_i c_i(z_i)
 \quad\text{for its inverse}.
\]
Its constant displacement is nilpotent.  No later substitution shifts
an arbitrary Laurent-coefficient series by a nonzero
$\sigma$-dependent constant, and no operation introduces relations
among the external parameters.  Artin approximation is applied only
after the final graph identities have been obtained, over the field
$K$ as in Lemma~\ref{lem:tensor-reduction}.

The central functions and the $G_i$ are still algebraic germs over $K$:
substitute the background values, which are elements of $K$, into their
finite algebraic presentations.  The selected background is critical
on every central edge and $G_i(h_i^0)=1$.  Moreover $F_i(h_i^0)=0$.
Indeed $\partial F_i$ vanishes along the formal background curve, so its
derivative with respect to $\sigma$ is zero; its value at $\sigma=0$
was zero.  We retain these normalizations after recentering.

The complementary tangent determinant and the gradient minors have
nonzero constant terms at the original joint origin.  Their evaluations
at the background are units of $\mathbb C[[\sigma]]$, and hence are
nonzero over $K$.  Thus the two recentered sections are still transverse
with complementary dimensions, and their gradient ideals still equal
the recentered residual ideals.  The latter may now have linear terms;
no later argument requires their original quadratic order.

\subsection{Conjugated laws on the full residual schemes}

The prescribed extra holonomies can be chosen explicitly.  In the basis
$(e,h,f)$ of $\mathfrak{sl}_2$, put
\[
 E=\ad(e)=\begin{pmatrix}0&-2&0\\0&0&1\\0&0&0\end{pmatrix},
 \qquad
 F=\ad(f)=\begin{pmatrix}0&0&0\\-1&0&0\\0&2&0\end{pmatrix}.
\]
For $a,b\in\{0,1,2\}$ prescribe
\begin{equation}
\label{eq:vertex-explicit-background}
 g_{ab}(\sigma)=\exp(a\sigma e)\exp(b\sigma f).
\end{equation}
The identity conjugator supplies $g_{00}$ and eight extra letters supply
the other eight matrices.  These choices can be made at both vertices.
They require only a finite list of conjugations in the geometric
construction, chosen before the cuts are made.

\begin{lemma}
\label{lem:vertex-adjoint-span}
The nine matrices $\Ad(g_{ab}(\sigma))$ span $\End_K(\g\otimes K)$.
\end{lemma}

\begin{proof}
Since $E^3=F^3=0$, the matrices are
\[
 (1+a\sigma E+(a\sigma)^2E^2/2)
 (1+b\sigma F+(b\sigma)^2F^2/2).
\]
Order $(a,b)$ lexicographically and vectorize matrices row by row.  The
determinant of the resulting nine columns is
\[
 512\sigma^{18}.
\]
For a direct calculation, expand in the nine ordered matrices
$E^pF^q/(p!q!)$, $0\leq p,q\leq2$.  Their vectorization determinant is
$8$.  The two evaluation matrices on $0,1,2$ are Vandermonde matrices
of determinant $2$, contributing $2^6$ to the tensor determinant.  The
sum of the exponents $p+q$ is $18$.  The displayed determinant is
nonzero in $K$.
\end{proof}

\begin{lemma}
\label{lem:vertex-triple-ideals}
For any three distinct active edges, the product of their three source
parameters belongs to $I_A$.  At $B$, every product of one matrix entry
from each of their three logarithmic $\beta$ holonomies belongs to
$I_B$.  The same conclusion holds with entries of $G_i-1$ in place of
the logarithmic holonomies.
\end{lemma}

\begin{proof}
Work in the recentered complete local ring over $K$ of one section,
again denoted by $R_P$, retaining the full square-free parameter
algebra at $A$.  Choose off-shell lifts
$X_i\in\g\otimes R_P$ of its marked meridian logarithms on
$R_P/I_P$, and write $\overline X_i$ for their quotient images.
At $A$ take $X_i=-t_iT_i$; at $B$ take the logarithmic holonomy of
the original vertex slice along $\beta_i$.  Each $X_i$ vanishes at
the background.  For three specified edges let $J_{ijk}$ be the ideal
in $R_P$ generated by all coordinate products containing one entry
from each of $X_i,X_j,X_k$.

Evaluate the marked boundary words at $\exp(X_i),\exp(X_j),\exp(X_k)$
and the fixed extra holonomies.  Modulo $I_P$ these are the actual
based holonomies of the flat vertex system, so the boundary laws of
Proposition~\ref{prop:geometric-reduction} make every scalar entry of a
law minus the identity lie in $I_P$.  The same entries lie in $J_{ijk}$:
the words are Brunnian in their three core meridians, so omitting any
one of the three makes the word the identity.  This latter assertion
is first an identity for independent formal matrix logarithms and is
then substituted into $R_P$; it does not require the chosen off-shell
connection to be flat.

Their cubic terms, up to the fixed signs and an invertible output
conjugation, are the scalar tests of
\begin{equation}
\label{eq:vertex-cubic-laws}
 [\Ad(g_a)X_i,[\Ad(g_b)X_j,\Ad(g_c)X_k]],
\end{equation}
with the three conjugators chosen independently.  The additional
core-word conjugations allowed by the geometric construction equal
the identity when the core logarithms vanish, and only change higher
terms.  In particular, their presence does not assert or require that
the globally based word is literally a fixed nested commutator.

By Lemma~\ref{lem:vertex-adjoint-span}, each input of the nonzero
trilinear double bracket can independently be acted on by the full
endomorphism algebra.  To see the spanning assertion directly, choose
a nonzero coefficient of this trilinear tensor and use rank-one
endomorphisms in its three inputs to select any three desired input
coordinates; choose an output functional detecting that coefficient.
Every scalar triple coordinate product is therefore a linear
combination of the cubic expressions in Equation~\eqref{eq:vertex-cubic-laws}.

All higher terms of a Brunnian law lie in $\mathfrak m J_{ijk}$, where
$\mathfrak m$ is the recentered maximal ideal, including the external
parameters.  If $Q$ is the ideal generated by the scalar law entries,
we have
\[
 Q\subset J_{ijk}\cap I_P,
 \qquad J_{ijk}=Q+\mathfrak m J_{ijk}.
\]
The module $J_{ijk}/Q$ is finitely generated over the complete
Noetherian local ring, so Nakayama's lemma gives $J_{ijk}=Q\subset I_P$.
This proves the assertion scheme-theoretically, including all
nilpotents and any syzygies in $r_P$.

At $A$ the quotient logarithms satisfy $\overline X_i=-t_iT_i$.
Choose a nonzero coordinate of each $T_i$ to deduce
$t_it_jt_k\in I_A$.  At $B$ use the corresponding $\beta$ logarithms.
On $B$'s flat quotient the image of $G_i(H_B)$ is
$\overline G_i=\exp(\overline X_i)$.  Every entry of
$\overline G_i-1$ is therefore in the ideal generated by the entries
of $\overline X_i$.  The triple products of the lifted entries
$G_i-1$ consequently belong to $I_B$ as well.  All earlier section
corrections were the
identity modulo the residual ideals, so the same conclusions hold on
the corrected sections.
\end{proof}

\subsection{Polynomial perturbations and the final graph families}
\label{subsec:vertex-polynomial}

The triple-ideal identities are now established on the full residual
schemes. The remaining task is to replace the logarithmic perturbation
by an algebraic function without losing exact vanishing or gradient
generation. Algebraic lifts of the holonomy matrices suffice: a
quadratic polynomial in those matrices has a controlled cubic error,
which can be removed on the source graph by a nilpotent correction.

For an active edge replace the logarithmic function by the polynomial
\begin{equation}
\label{eq:vertex-polynomial-weight}
 b_i(h_i)=-\epsilon_i\tr\left(
 T_i\left((G_i(h_i)-1)-\tfrac12(G_i(h_i)-1)^2\right)\right).
\end{equation}
These are algebraic germs over $K$, with zero constant term at the
recentered background.  We check the exact effects of this replacement.

On the plus section, setting $t_i=0$ gives, modulo $I_A$, a central
critical system on edge $i$ with both peripheral holonomies trivial:
$\alpha_i$ is constrained to be trivial and $\beta_i$ bounds in $K_A$.
The marked algebraic holonomy therefore satisfies
\begin{equation}
\label{eq:vertex-holonomy-ideal}
 G_i(H_A)-1\in (I_A,t_i)
\end{equation}
entry by entry.  This is equality in the quotient ring, not just
evaluation at its closed points.  The coordinate change used to absorb the difference in
Equation~\eqref{eq:vertex-log-mod-jacobian} is immaterial after setting $t_i=0$.

Let $R_i=\ell_i(G_i)-b_i$.  It has order at least three in the entries
of $G_i-1$, and each internal derivative has order at least two in
those entries.  Since $t_i^2=0$, Equation~\eqref{eq:vertex-holonomy-ideal} gives
\begin{align}
 t_iR_i(H_A)&\in t_i(I_A,t_i)^3=t_iI_A^3,\label{eq:vertex-taylor-value}\\
 t_i(\partial_hR_i)(H_A)&\in t_i(I_A,t_i)^2=t_iI_A^2.
                                                        \label{eq:vertex-taylor-gradient}
\end{align}
Thus the new total function on the existing plus section lies in
$\mathfrak t I_A^2$, where $\mathfrak t=(t_i:i\text{ active})$, and
its restricted gradient still has the form
$(M_A+\mathfrak t Q)r_A$ with a unit full-column minor.

The correction at the joint origin used the vanishing internal Hessian
and gained powers of the joint maximal ideal.  After recentering, the
Hessian may be nonzero.  Here the error instead contains a factor from
the nilpotent parameter ideal, and the following correction gains powers
of that ideal while keeping the same residual list.

\begin{lemma}[Correction over the full parameter algebra]
\label{lem:tensor-nilpotent-correction}
Let \(I\) be a finite set of cardinality \(m\), put
\(T=K[t_i:i\in I]/(t_i^2:i\in I)\), let \(J=(t_i:i\in I)\), and
work in \(T[[x]]\).  Fix a finite residual list \(r\), with ideal
\(\mathcal I=(r)\).
Suppose a formal section \(H\) and a formal potential \(S\) satisfy
\[
 S(H)\in J\mathcal I^2,\qquad
               (\partial_h S)(H)=M r,
\]
where \(M\) has a unit full-column minor.
Then a displacement of \(H\) in \(J\mathcal I\) makes \(S(H)=0\)
identically and keeps the restricted gradient ideal equal to
\(\mathcal I\).
If the original central parametrization has an injective tangent map,
the corrected section is still a smooth formal graph over \(T\).
\end{lemma}

\begin{proof}
Keep the functions \(r\) fixed on the parameter domain.
Suppose at one step the error is in \(J^q\mathcal I^2\), \(q\geq1\).
Write it as \(\sum_{a,b}c_{ab}r_a r_b\), with \(c_{ab}\in J^q\).
Using the inverse full-column minor of the restricted gradient matrix,
choose a displacement \(\Delta H\in J^q\mathcal I\) whose linear Taylor
term cancels this error.
The Taylor remainder belongs to
\[
       (\Delta H)^2\subset J^{2q}\mathcal I^2
                         \subset J^{q+1}\mathcal I^2 .
\]
This uses no vanishing of the internal Hessian.
The gradient changes by an element of \(J^q\mathcal I\), so its new
coefficient matrix is \(M+J^qQ\) for some matrix \(Q\).
Its chosen minor remains a unit because \(J\) is nilpotent.
Iterating terminates, since \(J^{m+1}=0\), and preserves gradient
generation throughout.
Every correction vanishes at \(t=0\).
A fixed linear projection that was invertible on the central tangent
therefore still has a unit relative Jacobian.  Formal inversion over
\(T\), fixing all parameters, gives the asserted graph.
The inverse reparametrization also pulls back every residual membership
witness.
\end{proof}

Apply Lemma~\ref{lem:tensor-nilpotent-correction} over $T[[y_A]]$, with
$J=\mathfrak t$, $\mathcal I=I_A$, and the fixed residual list $r_A$.
The preceding estimates give an error in
$\mathfrak t I_A^3\subset\mathfrak t I_A^2$ and a unit full-column
gradient minor against that same list.  The section is centered modulo
$\mathfrak t$: $H_A(0,0)=0$, so $H_A\in(y_A,\mathfrak t)^N$.
Formal substitution is continuous even with its nilpotent constants:
write $H_A=H_A(y_A,0)+\eta$, where $\eta\in\mathfrak t T[[y_A]]^N$;
expansion in $\eta$ is finite and substitution of $H_A(y_A,0)$ is
centered.  The central tangent remains injective by
Lemma~\ref{lem:vertex-linear} and the preceding coordinate changes.
Thus the lemma gives a corrected source section $H_L$ with exact
vanishing for the polynomial perturbations.  The corrections lie in
$\mathfrak t I_A$ and preserve the gradient ideal $I_A$.  The ideal
$\mathfrak t$ is nilpotent of exponent at most $m+1$; its square is not
set to zero.  In particular, mixed parameter monomials are retained at
every step.  A nonzero Hessian of a recentered central potential causes
no difficulty because the Taylor remainders gain powers of
$\mathfrak t$.

The minus function is still $S_0$ up to its harmless overall sign, so
the minus section needs no additional correction.  Denote it by $H_R$.
Every weight in Equation~\eqref{eq:vertex-polynomial-weight} has zero constant term
in the entries of $G_i-1$.  Lemma~\ref{lem:vertex-triple-ideals} therefore
gives $(b_ib_jb_k)(H_R)\in I_B$ for distinct active indices.  The source
parameter products remain in $I_A$.  Since the final restricted
gradients generate these same ideals, these are exactly
Equations~\eqref{eq:vertex-final-source} and~\eqref{eq:vertex-final-target}.

For completeness, the corrected source is a smooth graph over the
\emph{full} ring $T$, rather than just on its reduced central fiber.
Choose a linear projection $\rho$ on the internal coordinate space
which is invertible on the central tangent of the source.  The Jacobian
of $y_A\mapsto\rho H_L(y_A,t)$ is invertible modulo $(t)$ and hence is
a unit over $T[[y_A]]$.  The constant term of $\rho H_L$ can lie in the nilpotent ideal
$(t)$; translation by that constant is continuous for the
$(y_A,t)$-adic topology, which here agrees with the $y_A$-adic topology.
The formal inverse theorem over $T$ reparametrizes this map as
\[
 H_L(x,t)=(x,\Psi(x,t)),
\]
while leaving the parameters fixed.  All memberships pull back under
this same $T$-algebra automorphism.  At $t=0$ the last corrections
vanish, so complementarity with $H_R$ is preserved.  No flatness of
$T[[y_A]]/I_A$ over $T$ is required.

\begin{proof}[Proof of Proposition~\ref{prop:vertex-data}]
Use the smooth cut data and choose the finite extra-letter background
in Equation~\eqref{eq:vertex-explicit-background}.  Lemma~\ref{lem:vertex-linear},
the based connection slices, and Lemma~\ref{lem:vertex-offshell-identities}
give complementary formal sections, square-ideal errors, and gradient
generation by the actual residual lists.
Lemma~\ref{lem:vertex-correction} gives exact vanishing.  Apply the
separate marked algebraic edge-coordinate comparisons, absorb the
Jacobian terms, and recenter before extending to $K=\mathbb C((\sigma))$.
Lemma~\ref{lem:vertex-triple-ideals} supplies the exact source and target
triple-product memberships.  Finally
Equations~\eqref{eq:vertex-polynomial-weight}--\eqref{eq:vertex-taylor-gradient}
and the nilpotent correction produce algebraic $F_i,b_i$ and the
required final graphs.  The preceding graph reparametrization supplies
precisely the stated smooth families.  This proves all assertions.
\end{proof}

\begin{corollary}
\label{cor:vertex-tensor-interface}
The cut data of Proposition~\ref{prop:geometric-reduction}, with
$m\geq5$ and the indicated finite law list, supply all the hypotheses
of Theorem~\ref{thm:tensor-obstruction}.
\end{corollary}

\begin{proof}
Take the field, separate algebraic edge functions and graph families
of Proposition~\ref{prop:vertex-data}.  Equations
\eqref{eq:vertex-final-vanishing}--\eqref{eq:vertex-final-target} are
the exact equations and memberships required there.  The graphs may
mix the edge coordinates, while the fixed potentials remain separated
by edge.  Their external parameter ring is the complete square-free
algebra $T$, and their central tangent spaces are complementary.
Ordinary algebraic approximation and subsequent spreading out are
applied only at this final stage, as in the proof of
Theorem~\ref{thm:tensor-obstruction}; they are not applied to the earlier
arbitrary formal potentials or to a Laurent series term by term.
\end{proof}
